% Audit de veracite des citations de rapport de jury - 2026-09-20
% Bloc 4 : citation jury 2025 remplacee par le verbatim exact du
% Rapport_Jury_2025 (l.484) ; la version precedente en etait une reformulation.
% Citation jury 2024 (<< forme canonique >>) supprimee : zero occurrence de
% << forme canonique >> dans les 8 rapports 2018-2025, aucun remplacement pertinent.
\expandafter\ifx\csname CAPETINCLUDEMODE\endcsname\relax
\documentclass[11pt,a4paper]{article}
% preamble.tex — packages et commandes communes a toutes les fiches CAPET SII IE
\input{glyphtounicode}
\pdfgentounicode=1
\usepackage{cmap}
\usepackage[T1]{fontenc}
\usepackage[utf8]{inputenc}
\usepackage[french]{babel}
% babel[french] active a \begin{document} des raccourcis de ponctuation (:;!?)
% qui cassent la syntaxe to[...] de circuitikz (erreur "+ or - expected") ;
% on les desactive, ce qui ne change que l'espacement fin avant ces signes.
\AtBeginDocument{\shorthandoff{:;!?}}
\usepackage[a4paper, top=1.6cm, bottom=1.6cm, left=1.4cm, right=1.4cm]{geometry}
\setlength{\headheight}{14.5pt}
\usepackage{amsmath, amssymb, amsfonts}
\usepackage{siunitx}
\sisetup{locale=FR, output-decimal-marker={,}, inter-unit-product=\ensuremath{{}\cdot{}},
         group-minimum-digits=4}
% Unites non fournies par siunitx, declarees ici une fois pour toutes (fiches + exercices).
% Les gardes \ifdefined evitent tout conflit avec une redeclaration locale restee dans un
% fichier : ne pas retirer les gardes, ne pas redeclarer ces unites dans les fichiers.
\DeclareSIUnit{\jour}{j}
\ifdefined\var\else\DeclareSIUnit{\var}{var}\fi   % puissance reactive
\ifdefined\vad\else\DeclareSIUnit{\vad}{VAD}\fi   % puissance deformante
\ifdefined\tr\else\DeclareSIUnit{\tr}{tr}\fi      % tour (tr/min)
\usepackage{array, booktabs, multirow, tabularx}
\usepackage{xcolor}
\usepackage[most, breakable]{tcolorbox}
\usepackage{multicol}
\usepackage{enumitem}
\usepackage{mathtools}
\usepackage{lmodern}
% Refonte visuelle 2026 (direction "moderne coloré") : Lato pour le texte,
% maths sans-serif assorties via newtxsf (charge apres amsmath/amssymb).
\usepackage[default]{lato}
\usepackage{newtxsf}
\usepackage{microtype}
\usepackage{fancyhdr}
\usepackage{tikz}
\usetikzlibrary{positioning,arrows.meta}
\usepackage{circuitikz}
\usepackage{lastpage}

% ── Couleurs par domaine (palette modernisée 2026) ────────────────────
\definecolor{colElec}{RGB}{29, 78, 216}        % Électrotechnique — bleu vif
\definecolor{colElecLight}{RGB}{219, 234, 254} % fond bleu clair
\definecolor{colAuto}{RGB}{21, 128, 61}        % Automatique — vert
\definecolor{colMeca}{RGB}{194, 65, 12}        % Mécanique — orange
\definecolor{colTherm}{RGB}{185, 28, 28}       % Thermique — rouge
\definecolor{colQT}{RGB}{232, 244, 233}        % (hérité, plus utilisé par questiontype)
\definecolor{colJury}{RGB}{255, 248, 225}      % fond jury (crème)
\definecolor{colJuryFrame}{RGB}{202, 138, 4}   % bordure/bandeau jury (ambre)
\definecolor{colGrayLight}{RGB}{245, 245, 247} % fond gris neutre
\definecolor{colSeparator}{RGB}{180, 180, 180} % gris séparateur
\definecolor{colExo}{RGB}{8, 145, 178}          % Banque d'exercices — cyan (distinct des fiches)
% ── Encre et gris de la refonte ───────────────────────────────────────
\definecolor{colInk}{RGB}{23, 32, 42}          % texte courant (quasi-noir)
\definecolor{colSlate}{RGB}{83, 98, 118}       % gris ardoise (méta, en-têtes)
\definecolor{colRule}{RGB}{203, 213, 225}      % filets clairs
\definecolor{colJuryBg}{RGB}{255, 248, 225}    % fond crème jury
\definecolor{colPresence}{RGB}{83, 98, 118}    % bloc 5 (présence dans les sujets)
\AtBeginDocument{\color{colInk}}

% ── Style encadré exercice + corrigé (banque d'exercices) ─────────────
\tcbset{
  exoBox/.style={
    breakable,
    enhanced,
    colback=colExo!8!white,
    colframe=colExo,
    colbacktitle=colExo,
    coltitle=white,
    fonttitle=\bfseries\small,
    left=4pt, right=4pt, top=3pt, bottom=3pt,
    before skip=5pt, after skip=5pt,
  }
}

% ── Bascule enonce / corrige (banque d'exercices) ─────────────────────
% \CAPETCORRIGE est defini en ligne de commande (pdflatex "\def\CAPETCORRIGE{1}\input{...}")
% pour compiler la variante corriges d'un meme fichier source ; par defaut (non defini),
% la compilation produit la variante enonces. \CAPETSHOW{enonce}{corrige} choisit la branche.
%
% TROIS MODES (refonte 2026-09-20, manuel de revision) :
%   - aucun drapeau   -> imprime l'enonce            (variante enonces, INCHANGE)
%   - \CAPETCORRIGE   -> imprime le corrige          (variante corriges, INCHANGE)
%   - \CAPETMANUEL    -> imprime l'enonce ET met le corrige de cote pour l'annexe
%
% Le mode manuel stocke chaque corrige dans une macro (\csgdef), sans ecriture
% fichier ni re-tokenisation : le contenu (circuitikz, accents, %) traverse intact.
% L'annexe les rejoue ensuite dans l'ordre via \CAPETANNEXE.
\usepackage{etoolbox}

\newcounter{CAPETexo}          % numero continu d'exercice dans le manuel
\setcounter{CAPETexo}{0}
\newcounter{CAPETannexe}       % compteur de relecture en annexe

% Titre courant de la notion en cours (pose par \CAPETNOTION), sert d'intertitre annexe
\newcommand{\CAPETnotioncode}{}
\newcommand{\CAPETnotiontitre}{}

% \CAPETNOTION{code}{titre} : ouvre une notion dans le manuel (memorise le libelle,
% et marque le rang du premier exercice de la notion pour l'intertitre d'annexe).
\newcommand{\CAPETNOTION}[2]{%
  \renewcommand{\CAPETnotioncode}{#1}%
  \renewcommand{\CAPETnotiontitre}{#2}%
  \csgdef{CAPETnotionstart\the\numexpr\value{CAPETexo}+1\relax}{#1 -- #2}%
}

% ── \degree robuste ───────────────────────────────────────────────────
% Certains fichiers d'exercices utilisent \degree. Ni le preambule ni siunitx
% ne le definissent, et \providecommand local ne suffit pas : sous hyperref,
% \degree existe deja comme commande factice (\ERROR), donc \providecommand
% s'abstient et l'appel echoue a la compilation du manuel (erreur invisible
% en compilation isolee). On le (re)definit donc inconditionnellement ici,
% apres le chargement des packages, pour tout le corpus.
\AtBeginDocument{\providecommand{\degree}{}\renewcommand{\degree}{\ensuremath{^\circ}}}

% Libelle decoratif : les fichiers d'exercices utilisent aussi \CAPETSHOW dans
% leur \fancyhead et leur bloc de titre, pour afficher "(enonces)" ou "(corriges)".
% Ces appels-la ne sont PAS des exercices : en mode manuel ils ne doivent ni
% consommer un numero, ni alimenter l'annexe. On les reconnait a leur premier
% argument, strictement egal au libelle ci-dessous.
\newcommand{\CAPETLIBELLEENONCES}{(\'enonc\'es)}

\ifdefined\CAPETMANUEL
  % --- Mode manuel : enonce en place, corrige differe vers l'annexe ---
  \newcommand{\CAPETSHOW}[2]{%
    \def\CAPETargun{#1}%
    \ifx\CAPETargun\CAPETLIBELLEENONCES
      #1%            appel decoratif : rendu tel quel, hors numerotation
    \else
      \CAPETSHOWEXO{#1}{#2}%
    \fi
  }
  \newcommand{\CAPETSHOWEXO}[2]{%
    \stepcounter{CAPETexo}%
    \csgdef{CAPETcorrige\theCAPETexo}{#2}%
    \edef\CAPETtmpnum{\theCAPETexo}%
    #1%
    \par\vspace{2pt}%
    {\footnotesize\itshape\color{colExo}%
     \hyperref[capet:cor:\CAPETtmpnum]{Corrige de l'exercice \CAPETtmpnum\ $\rightarrow$ p.~\pageref{capet:cor:\CAPETtmpnum}}%
     \label{capet:eno:\CAPETtmpnum}}%
  }
  % \CAPETANNEXE : rejoue tous les corriges captures, dans l'ordre, avec intertitres
  \newcommand{\CAPETANNEXE}{%
    \setcounter{CAPETannexe}{0}%
    \loop\ifnum\value{CAPETannexe}<\value{CAPETexo}%
      \stepcounter{CAPETannexe}%
      \edef\CAPETcurnum{\theCAPETannexe}%
      \ifcsname CAPETnotionstart\CAPETcurnum\endcsname
        \par\medskip
        {\normalsize\bfseries\color{colExo}\csuse{CAPETnotionstart\CAPETcurnum}}%
        \par\nobreak\smallskip
      \fi
      \begin{tcolorbox}[exoBox, breakable,
        title={Corrige de l'exercice \CAPETcurnum\
               \hfill\normalfont\footnotesize
               \hyperref[capet:eno:\CAPETcurnum]{enonce p.~\pageref{capet:eno:\CAPETcurnum}}}]%
      \label{capet:cor:\CAPETcurnum}%
      \csuse{CAPETcorrige\CAPETcurnum}%
      \end{tcolorbox}%
    \repeat
  }
\else
  % --- Modes historiques, comportement strictement inchange ---
  \newcommand{\CAPETSHOW}[2]{\ifdefined\CAPETCORRIGE #2\else #1\fi}
  \newcommand{\CAPETANNEXE}{}
\fi

% ── Difficulté 1 à 3 (étoiles pleines, amssymb) ───────────────────────
\newcommand{\difficulte}[1]{%
  \ifcase#1\relax
  \or $\bigstar$\;niveau 1/3%
  \or $\bigstar\bigstar$\;niveau 2/3%
  \or $\bigstar\bigstar\bigstar$\;niveau 3/3%
  \fi
}

% ── Style encadré extrait de document technique (banque d'exercices) ──
\tcbset{
  docTech/.style={
    breakable,
    enhanced,
    colback=colGrayLight,
    colframe=colSeparator,
    fonttitle=\bfseries\small\itshape,
    coltitle=black,
    colbacktitle=colGrayLight,
    left=4pt, right=4pt, top=3pt, bottom=3pt,
    before skip=4pt, after skip=4pt,
  }
}

% ── Style tcolorbox de section principale ────────────────────────────
% Refonte 2026 : bandeau titre coloré arrondi + fond très pâle, sans cadre.
\tcbset{
  ficheSection/.style={
    breakable, enhanced jigsaw,
    frame hidden, boxrule=0pt, arc=1.8mm,
    colback=#1!4!white,
    colbacktitle=#1,
    coltitle=white,
    fonttitle=\sffamily\bfseries,
    left=7pt, right=7pt, top=5pt, bottom=5pt,
    toptitle=3pt, bottomtitle=3pt,
    before skip=10pt, after skip=8pt,
  }
}

% ── Style encadré question type ──────────────────────────────────────
% Carte blanche, bandeau sombre, liseré très clair.
\tcbset{
  questiontype/.style={
    breakable, enhanced jigsaw,
    colback=white,
    frame hidden, boxrule=0pt, arc=1.8mm,
    colbacktitle=colElec!78!black,
    coltitle=white,
    fonttitle=\sffamily\bfseries\small,
    left=7pt, right=7pt, top=4pt, bottom=4pt,
    toptitle=2.5pt, bottomtitle=2.5pt,
    before skip=8pt, after skip=8pt,
    borderline={0.6pt}{0pt}{colElec!12!white},
  }
}

% ── Style encadré jury ───────────────────────────────────────────────
% Bandeau ambre, fond crème, sans cadre.
\tcbset{
  jury/.style={
    breakable, enhanced jigsaw,
    colback=colJuryBg,
    frame hidden, boxrule=0pt, arc=1.8mm,
    colbacktitle=colJuryFrame,
    coltitle=white,
    fonttitle=\sffamily\bfseries,
    left=7pt, right=7pt, top=4pt, bottom=4pt,
    toptitle=3pt, bottomtitle=3pt,
    before skip=10pt, after skip=8pt,
  }
}

% ── En-tête / pied de page communs (refonte 2026) ─────────────────────
% #1 code fiche, #2 titre court, #3 couleur domaine
\newcommand{\ficheHeaders}[3]{%
  \fancyhead[L]{\small\sffamily{\bfseries\textcolor{#3}{#1}}\hspace{0.7em}\textcolor{colSlate}{#2}}%
  \fancyhead[R]{\small\sffamily\textcolor{colSlate}{CAPET SII IE\, --- \,Session 2026}}%
  \fancyfoot[C]{\small\sffamily\textcolor{colSlate}{\thepage\,/\,\pageref{LastPage}}}%
  \fancyhead[C]{}\fancyfoot[L]{}\fancyfoot[R]{}%
}

% ── Bloc en-tête de fiche (BLOC 0, refonte 2026) ──────────────────────
% #1 code  #2 titre  #3 couleur domaine  #4 domaine  #5 ligne présence  #6 sous-thèmes
\newcommand{\ficheEntete}[6]{%
  \begin{tcolorbox}[enhanced, breakable,
    frame hidden, boxrule=0pt, arc=2.5mm,
    interior style={left color=#3, right color=#3!75!black},
    left=9pt, right=9pt, top=7pt, bottom=7pt,
    before skip=0pt, after skip=6pt]
    {\color{white}\sffamily\bfseries\LARGE #1\quad #2}\\[4pt]
    {\color{#3!22!white}\sffamily\small \MakeUppercase{#4}\quad\textbar\quad #5}\\[1pt]
    {\color{#3!35!white}\sffamily\footnotesize Sous-thèmes : #6}
  \end{tcolorbox}%
}

% ── Mise en page ─────────────────────────────────────────────────────
\pagestyle{fancy}
\fancyhf{}
\renewcommand{\headrulewidth}{0.5pt}
\renewcommand{\footrulewidth}{0pt}
\AtBeginDocument{%
  \renewcommand{\headrule}{{\color{colRule}\hrule width\headwidth height 0.5pt}}%
}

\setlist[itemize]{leftmargin=1.2em, itemsep=1pt, topsep=2pt, parsep=0pt}
\setlist[enumerate]{leftmargin=1.2em, itemsep=3pt, topsep=2pt, parsep=0pt}

\setlength{\parskip}{2pt}
\setlength{\columnsep}{10pt}

\def\CAPETBEGINFILE{\begin{document}}
\def\CAPETENDFILE{\end{document}}
\else
\def\CAPETBEGINFILE{}
\def\CAPETENDFILE{}
\fi

\ficheHeaders{T26}{Modélisation des systèmes}{colAuto}

\CAPETBEGINFILE

% -----------------------------------------------------------------------
%  BLOC EN-TETE
% -----------------------------------------------------------------------

\ficheEntete{T26}{Modélisation des systèmes linéaires (automatique)}{colAuto}%
  {Automatique - Modélisation et analyse des systèmes}%
  {\textbf{Présence :} \textbf{9/13} sujets (2013-2026) - thème \emph{fondamental}}%
  {transformée de Laplace - FdT 1\textsuperscript{er}/2\textsuperscript{e} ordre - identification graphique - schéma-bloc FTBF - modèle MCC - diagramme de Bode - représentation d'état}

% -----------------------------------------------------------------------
%  1. RAPPELS DE COURS
% -----------------------------------------------------------------------

\begin{tcolorbox}[
  ficheSection=colAuto,
  title={1. Rappels de cours},
]

% -- 1.1 Transformee de Laplace ----------------------------------------
\textbf{1.1 Transformée de Laplace - transformées usuelles}

\smallskip
\begin{center}
\renewcommand{\arraystretch}{1.3}
\begin{tabular}{lll}
\toprule
\textbf{$f(t)$} & \textbf{$F(p)$} & \textbf{Remarque} \\
\midrule
Échelon $u(t)$ & $1/p$ & \\
Rampe $t\,u(t)$ & $1/p^2$ & \\
Exponentielle $e^{-t/\tau}$ & $\dfrac{1}{1+\tau p}$ & \\
Dérivée $f'(t)$ & $p\,F(p)$ & CI nulles (Heaviside) \\
Intégrale $\int_0^t f$ & $F(p)/p$ & \\
\bottomrule
\end{tabular}
\end{center}

\medskip\hrule\vspace{4pt}

% -- 1.2 FdT 1er ordre ------------------------------------------------
\textbf{1.2 Système du 1\textsuperscript{er} ordre - forme canonique}

\[
  H(p) = \frac{K}{1 + \tau\,p}
\]
Réponse indicielle : $s(t) = K\,E_0\,(1 - e^{-t/\tau})$ \quad;\quad $t_{r5\%} \approx 3\tau$

\smallskip
\begin{center}
\begin{tikzpicture}[scale=0.9, font=\footnotesize]
  \draw[->] (0,0) -- (6.3,0) node[right] {$t$};
  \draw[->] (0,-0.2) -- (0,2.4) node[above] {$s(t)$};
  \draw[thick, colAuto, domain=0:6, samples=60, smooth] plot (\x, {2*(1-exp(-\x/1.5))});
  \draw[dashed] (0,2) -- (6.2,2);
  \node[right] at (6.2,2) {\scriptsize $K E_0$};
  \draw[dashed] (0,1.264) -- (1.5,1.264) -- (1.5,0);
  \node[below] at (1.5,0) {\scriptsize $\tau$};
  \node[left] at (0,1.264) {\scriptsize $0{,}632\,KE_0$};
\end{tikzpicture}
\end{center}

\textbf{Identification sur l'essai indiciel :}
\begin{itemize}
  \item $K = s(\infty)/E_0$ (valeur finale / échelon)
  \item $\tau$ : temps pour atteindre $63{,}2\,\%$ de $K\,E_0$ (méthode tangente)
\end{itemize}

\medskip\hrule\vspace{4pt}

% -- 1.3 FdT 2e ordre -------------------------------------------------
\textbf{1.3 Système du 2\textsuperscript{e} ordre - forme canonique}

\[
  H(p) = \frac{K}{1 + \dfrac{2m}{\omega_0}\,p + \dfrac{p^2}{\omega_0^2}}
  \qquad
  D\% = 100\,e^{-m\pi/\sqrt{1-m^2}} \quad (m<1)
\]
\begin{center}
\renewcommand{\arraystretch}{1.2}
\begin{tabular}{lll}
\toprule
$m > 1$ & Sur-amorti & Pas d'oscillations \\
$m = 1$ & Critique & Retour le plus rapide sans dépassement \\
$m < 1$ & Sous-amorti & Oscillations et résonance \\
\bottomrule
\end{tabular}
\end{center}

\medskip\hrule\vspace{4pt}

% -- 1.4 Modele de la MCC ---------------------------------------------
\textbf{1.4 Modèle de la machine à courant continu (MCC)}

\[
  u(t) = R\,i(t) + L\,\frac{di}{dt} + e(t)
  \qquad;\qquad
  e(t) = K_e\,\Omega(t)
  \qquad;\qquad
  K_e = K_c
\]

En régime permanent : $U = R\,I + K_e\,\Omega$ ; fonction de transfert moteur-mécanisme :
\[
  \frac{\Omega(p)}{U(p)} = \frac{K_m}{1 + \tau_m\,p}
  \qquad K_m = \frac{1}{K_e} \qquad \tau_m = \frac{J\,R}{K_e^2 + f_v\,R}
\]

\medskip\hrule\vspace{4pt}

% -- 1.5 Schema-bloc - boucle fermee ----------------------------------
\textbf{1.5 Schéma-bloc - règles de manipulation}

\[
  H_{BF}(p) = \frac{H_D(p)}{1 + H_D(p)\,H_R(p)} \quad \text{(retour négatif)}
  \qquad
  H_{BO}(p) = H_D(p)\,H_R(p)
\]

Erreur statique en poursuite (classe 0) :
\[
  \varepsilon_s = \frac{E}{1 + K_{BO}} \qquad (\text{nulle si intégrateur en BO})
\]

\medskip\hrule\vspace{4pt}

% -- 1.6 Representation d'etat -----------------------------------------
\textbf{1.6 Représentation d'état}

\[
  \dot{\mathbf{x}}(t) = A\,\mathbf{x}(t) + B\,\mathbf{u}(t)
  \qquad;\qquad
  \mathbf{y}(t) = C\,\mathbf{x}(t) + D\,\mathbf{u}(t)
\]
Utilisée pour le modèle moyen d'un convertisseur élévateur (2025 E1).

\end{tcolorbox}

% -----------------------------------------------------------------------
%  2. FORMULAIRE SYNTHETIQUE
% -----------------------------------------------------------------------

\begin{tcolorbox}[
  ficheSection=colAuto,
  title={2. Formulaire synthétique},
]
\small
\begin{multicols}{2}
\renewcommand{\arraystretch}{1.45}
\begin{tabular}{@{}p{5cm}ll@{}}
\toprule
\textbf{Formule} & \textbf{Usage} & \textbf{U.} \\
\midrule
\multicolumn{3}{@{}l}{\footnotesize\textit{Transformée de Laplace}} \\
$\mathcal{L}\{f'(t)\} = p\,F(p)$ & CI nulles & - \\
$\mathcal{L}\{e^{-t/\tau}\} = 1/(1+\tau p)$ & Expo & - \\
\midrule
\multicolumn{3}{@{}l}{\footnotesize\textit{FdT 1\textsuperscript{er} ordre}} \\
$H(p) = K/(1+\tau p)$ & Forme canon. & - \\
$t_{r5\%} \approx 3\tau$ & Rapidité & s \\
$K = s(\infty)/E_0$ & Identification & - \\
$\tau$ : à $63{,}2\,\%$ de $K\,E_0$ & Identification & s \\
\bottomrule
\end{tabular}

\columnbreak

\renewcommand{\arraystretch}{1.45}
\begin{tabular}{@{}p{5cm}ll@{}}
\toprule
\textbf{Formule} & \textbf{Usage} & \textbf{U.} \\
\midrule
\multicolumn{3}{@{}l}{\footnotesize\textit{FdT 2\textsuperscript{e} ordre}} \\
$H = K/(1+2mp/\omega_0+p^2/\omega_0^2)$ & Forme canon. & - \\
$D\% = 100\,e^{-m\pi/\sqrt{1-m^2}}$ & Dépassement & \% \\
\midrule
\multicolumn{3}{@{}l}{\footnotesize\textit{Boucle fermée}} \\
$H_{BF} = H_D/(1+H_{BO})$ & FTBF & - \\
$\varepsilon_s = E/(1+K_{BO})$ & Erreur statique & - \\
\midrule
\multicolumn{3}{@{}l}{\footnotesize\textit{MCC}} \\
$U = RI + K_e\Omega$ & Régime perm. & V \\
$K_e = K_c$ & [V\,s/rad] = [N\,m/A] & \\
\bottomrule
\end{tabular}
\end{multicols}
\normalsize
\end{tcolorbox}

% -----------------------------------------------------------------------
%  3. QUESTIONS TYPES
% -----------------------------------------------------------------------

\begin{tcolorbox}[
  ficheSection=colAuto,
  title={3. Questions types - extraites des sujets réels},
]

% -- QT1 -------------------------------------------------------------

\begin{tcolorbox}[questiontype,
  title={QT1 - MCC - équation électrique et identification de $K_e$ \hfill\normalfont\footnotesize E1 2026 (A1-A4)}
]
\textbf{Énoncé.}
MCC d'une porte de bus ($R = 6{,}6\,\Omega$ ; $L = 100\,\mathrm{mH}$ ; $K_e = 1{,}044\,\mathrm{V\!\cdot\!s\!\cdot\!rad^{-1}}$). Écrire l'équation électrique de l'induit. Simplifier en régime permanent. Exprimer $R$ à partir d'un essai rotor bloqué ($\Omega = 0$). Retrouver $K_e$ à partir des conditions nominales.

\tcbline

\textbf{Corrigé.}

\textbf{Étape 1 — Équation électrique de l'induit} (loi des mailles) :
\[
  u(t) = R\,i(t) + L\,\frac{di}{dt} + e(t) \quad;\quad e(t) = K_e\,\Omega(t)
\]

\textbf{Étape 2 — Régime permanent} ($di/dt = 0$) : $U = R\,I + K_e\,\Omega$.

\textbf{Étape 3 — Essai rotor bloqué} ($\Omega = 0$, $e = 0$) : $U_{bl} = R\,I_{bl}$ - $R = U_{bl}/I_{bl}$.

\textbf{Étape 4 — $K_e$ à partir du nominal.} $K_e = (U_{nom} - R\,I_{nom})/\Omega_{nom}$.

\textbf{Conclusion.} Ces deux essais complémentaires (rotor bloqué pour $R$, point nominal pour $K_e$) permettent d'identifier complètement le modèle électrique de la MCC à partir de mesures accessibles sans capteur de couple.
\end{tcolorbox}

\vspace{4pt}

% -- QT2 -------------------------------------------------------------

\begin{tcolorbox}[questiontype,
  title={QT2 - FdT depuis l'essai indiciel - 1\textsuperscript{er} ordre \hfill\normalfont\footnotesize E1 2026 (C4)}
]
\textbf{Énoncé.}
L'essai indiciel d'un élément de la chaîne d'asservissement montre une réponse qui part de 0 et atteint une valeur finale. Proposer un ordre. Construire graphiquement $K$ et $\tau$. Exprimer $H_2(p)$ sous forme canonique.

\tcbline

\textbf{Corrigé.}

\textbf{Étape 1 — Ordre du système.} La réponse exponentielle croissante sans oscillation - système du \textbf{1\textsuperscript{er} ordre}.

\textbf{Étape 2 — Lecture graphique.}
\begin{itemize}
  \item $K = s(\infty)/E_0$ : lire la valeur finale divisée par l'échelon.
  \item $\tau$ : abscisse du point à $63{,}2\,\%$ de $s(\infty)$ (ou intersection de la tangente à l'origine avec le plateau).
\end{itemize}

\textbf{Étape 3 — Forme canonique.} $H_2(p) = \dfrac{K}{1 + \tau\,p}$ avec les valeurs lues numériquement.

\textbf{Conclusion.} Cette identification purement graphique, sans connaître les paramètres physiques internes de l'élément, suffit à obtenir un modèle exploitable pour la suite de l'étude de la boucle d'asservissement.
\end{tcolorbox}

\vspace{4pt}

% -- QT3 -------------------------------------------------------------

\begin{tcolorbox}[questiontype,
  title={QT3 - FTBF du 2\textsuperscript{e} ordre - dépassement nul \hfill\normalfont\footnotesize E1 2026 (C7-C8)}
]
\textbf{Énoncé.}
La FTBF d'un asservissement avec correcteur $K_C$ est :
$H_{BF}(p) = \dfrac{K_C\,K_m}{1 + (R + K_C\,K_m\,K_{gt})/(K_e)\,\tau_m\,p + \tau_m\,p^2}$.
Montrer qu'elle se met sous la forme canonique du 2\textsuperscript{e} ordre. Exprimer $m$ en fonction de $K_C$. En déduire $K_C$ assurant $D\% = 0$.

\tcbline

\textbf{Corrigé.}

\textbf{Étape 1 — Identification par lecture de la forme canonique.}
\[
  \omega_0^2 = \frac{K_C\,K_m}{\tau_m} \quad;\quad 2m\,\omega_0 = \frac{R + K_C\,K_m\,K_{gt}}{K_e\,\tau_m}
\]

\textbf{Étape 2 — Condition de dépassement nul.} $D\% = 0 \Leftrightarrow m \geq 1$. Résoudre $m(K_C) = 1$ - valeur limite $K_{Clim}$.

\textbf{Conclusion.} Pour $K_C \leq K_{Clim}$, le dépassement est annulé ($m \geq 1$) : c'est cette borne supérieure sur $K_C$, et non une valeur unique, qui définit la plage de réglages acceptables — il reste ensuite à vérifier que la marge de phase $M_\varphi$ (Bode) est également respectée sur cette plage.
\end{tcolorbox}

\end{tcolorbox}

% -----------------------------------------------------------------------
%  4. ATTENTES DU JURY
% -----------------------------------------------------------------------

\begin{tcolorbox}[
  jury,
  title={4. Attentes du jury - Erreurs récurrentes et points de vigilance},
]

\begin{enumerate}[label=\textbf{\arabic*.}, leftmargin=*, itemsep=5pt]

  \item \textbf{Maîtriser les conditions de Heaviside.} Tout passage équation différentielle → FdT suppose des conditions initiales nulles. Le mentionner explicitement avant d'appliquer la transformée de Laplace.

  \item \textbf{Mettre sous forme canonique avec démonstration.} Recopier la formule finale sans démonstration est sanctionné. Dérouler chaque étape d'identification.

  \item \textbf{Identification à $63{,}2\,\%$, pas à $50\,\%$ ni à $90\,\%$.} L'erreur sur la lecture de $\tau$ est fréquente - tracer la droite horizontale à $0{,}632 \times K\,E_0$ sur le document réponse.

  \item \textbf{Homogénéité et unités.} $K_e$ [V\,s\,rad$^{-1}$] = $K_c$ [N\,m\,A$^{-1}$] numériquement - ne pas les confondre. Les unités du gain statique $K$ dépendent de celles de sortie/entrée.

  \item \textbf{Signe du retour dans $H_{BF}$.} $H_{BF} = H_D/(1 + H_{BO})$ : le signe $+$ au dénominateur correspond à un retour \textbf{négatif}. Un retour positif donnerait $1 - H_{BO}$ - instabilité.

  \item \textbf{Thème le plus décrié.} Jury 2023 : "L'automatique est un champ disciplinaire auquel les candidats ne peuvent se soustraire. Le jury invite les candidats à revoir les concepts de base." - réviser en priorité.

\end{enumerate}

\medskip
\begin{tcolorbox}[
  colback=colJuryFrame!12!colJury,
  colframe=colJuryFrame,
  left=4pt, right=4pt, top=2pt, bottom=2pt,
  fontupper=\small\itshape
]
Rapport jury 2025 : «~De nombreux candidats ne maitrisent pas les lois basiques de l'électricité (la loi des nœuds, la loi des mailles, les calculs de puissance ou encore la loi d'Ohm).~»
\end{tcolorbox}

\end{tcolorbox}

% -----------------------------------------------------------------------
%  PIED - PRESENCE DANS LES SUJETS
% -----------------------------------------------------------------------

\begin{tcolorbox}[
  ficheSection=colPresence,
  title={Présence dans les sujets analysés (2013-2026)},
]
\begin{center}
\footnotesize\renewcommand{\arraystretch}{1.15}
\begin{tabular}{lccccccccccccc}
\toprule
\textbf{Année} & 13 & 14 & 15 & 16 & 17 & 18 & 19 & 20 & 21 & 22 & 23 & 24 & 25 \\
\midrule
E1 (commune)   & - & - & - & $\checkmark$ & - & $\checkmark$ & $\checkmark$ & - & $\checkmark$ & - & $\checkmark$ & - & $\checkmark$ \\
E2 discipl. IE & - & $\checkmark$ & - & - & - & - & - & $\checkmark$ & - & - & $\checkmark$ & $\checkmark$ & $\checkmark$ \\
E2 appl. comm. & - & - & - & - & - & - & - & - & - & - & - & $\checkmark$ & - \\
\bottomrule
\end{tabular}
\end{center}
\smallskip
\footnotesize\textbf{Sous-thèmes fréquents :}
équation électrique MCC + régime permanent - FdT 1\textsuperscript{er} ordre (identification graphique $K$, $\tau$) - FTBF 2\textsuperscript{e} ordre ($m$, $\omega_0$) - dépassement nul - erreur statique (classe 0 vs 1)

\medskip
\textbf{Ordres de grandeur (repères rapides).}
\begin{itemize}[leftmargin=1.3em]
  \item Réseau BT FR : tension phase-neutre typique \SI{230}{\volt}, tension phase-phase typique \SI{400}{\volt}, fréquence \SI{50}{\hertz}.
  \item Courants d'étude fréquents en sujet : de quelques \si{\ampere} à quelques dizaines de \si{\ampere}.
  \item En conversion d'énergie, un rendement attendu reste généralement entre 0{,}8 et 0{,}98 selon la technologie.
\end{itemize}

\medskip
\textbf{Synthèse express (à refaire sans support).}
\begin{enumerate}[leftmargin=1.4em]
  \item Poser les hypothèses du modèle et les conventions de signe.
  \item Écrire 2 à 4 relations clés avec unités.
  \item Réaliser un calcul numérique complet et vérifier l'ordre de grandeur.
  \item Justifier le résultat par un argument technique de niveau CAPET.
\end{enumerate}

\end{tcolorbox}

\CAPETENDFILE
